\documentclass[english, parskip=full]{scrartcl}
\usepackage[utf8]{inputenc}  % Kodierung der Datei
\usepackage[english]{babel}  % Sprache auf Deutsch 
\usepackage{color}
\usepackage{graphicx, gensymb}
\usepackage{amsfonts, cancel, amssymb}
\usepackage{amsmath, amsthm, array, esint}
\usepackage{booktabs, multirow}  % schönere Tabellen
\usepackage{caption}  % Anpassung der Captions
\usepackage{scrlayer-scrpage}    % Manipulation des Seitenstils
\usepackage[hidelinks]{hyperref}  % Links und weitere PDF-Features
\usepackage{typearea, tikz, pgffor, verbatim}
\usetikzlibrary{calc, quotes}
\usepackage[cache=false, outputdir=build]{minted}
\usepackage{placeins} % provides \FloatBarrier

\definecolor{bg}{rgb}{0.9,0.9,0.9}
\newcommand\numberthis{\addtocounter{equation}{1}\tag{\theequation}}
\newcommand{\kom}{}
\renewcommand{\i}{\text{i}}
\newcommand{\e}{\text{e}}
\newcommand*{\Fourier}{\textsc{Fourier}\ }
\newcommand*{\F}{\mathcal{F}}
\newcommand*{\sinc}{\text{sinc\,}}
\newcommand{\conv}{\mathop{\scalebox{3.5}{\raisebox{-0.38ex}{\makebox[0.5\width][c]{$\ast$}}}}\limits}

\newcommand*{\titel}{01 Das Gauss-Seidel-Verfahren}
\newcommand*{\autor}{Joachim Schwardt und Julian Fleck}

\hypersetup{pdfauthor={\autor}, pdftitle={\titel}}

\subject{Numerik II - Aufgaben}
\title{\titel}
\author{\begin{tabular}{l|l}
Joachim Schwardt & 4768711 \\ 
Julian Fleck & 4759587\\ 
\end{tabular}}
\date{\today}

\begin{document}

\maketitle  % Titel setzen

\section*{Aufgabe 1}
Gegeben ist das lineare Gleichungssystem $Ax = b$ mit $b=(9,6,4)^\top$, $x=(x_1, x_2, x_3)^\top$ und
\begin{align}
A &= \begin{pmatrix}
8&2&1 \\ 2&6&2 \\ 2&1&4
\end{pmatrix}
\end{align}
mit dem Startwert $x^0 = (0,0,0)^\top$.
Die Iterationsvorschrift des Gauss-Seidel-Verfahrens lautet hier explizit
\begin{align}
x^{k+1}_1 &= \frac{1}{a_{11}} \left( b_1 - a_{12}x^k_2 - a_{13}x^k_3 \right),\\
x^{k+1}_2 &= \frac{1}{a_{22}} \left( b_2 - a_{21}x^{k+1}_1 - a_{23}x^k_3 \right),\\
x^{k+1}_3 &= \frac{1}{a_{33}} \left( b_3 - a_{31}x^{k+1}_1 - a_{32}x^{k+1}_2 \right).
\end{align}
Für den gegebenen Startwert erhält man damit
\begin{align}
x^1_1 &= \frac{1}{8} \left( 9 \right) = \frac{9}{8},\\
x^1_2 &= \frac{1}{6} \left( 6 - 2\cdot\frac{9}{8} \right) = \frac{5}{8},\\
x^1_3 &= \frac{1}{4} \left( 4 - 2\cdot\frac{9}{8} - 1\cdot\frac{5}{8} \right) = \frac{9}{32}.
\end{align}
Der zweite Iterationsschritt lautet
\begin{align}
x^2_1 &= \frac{1}{8} \left( 9 - 2\cdot \frac{5}{8} - 1\cdot\frac{9}{32} \right) = \frac{239}{256},\\
x^2_2 &= \frac{1}{6} \left( 6 - 2\cdot \frac{239}{256} - 2\cdot \frac{9}{32} \right) = \frac{457}{768},\\
x^2_3 &= \frac{1}{4} \left( 4 - 2\cdot \frac{239}{256} - 1\cdot \frac{457}{768} \right) =  \frac{1181}{3072}.
\end{align}




\section*{Aufgabe 2}
Sei $T_\text{GS} := I - CA$ die Iterationsmatrix des Gauss-Seidel-Verfahrens.
Dann gilt 
\begin{align}
(C^{-1})_{jk} &= \begin{cases} a_{jk} & j \ge k \\ 0 & j < k \end{cases}.
\end{align}
Insbesondere gilt dann für die Matrix $M := C^{-1} - A$ 
\begin{align}
\det(M) &= 0,
\end{align}
da $(M)_{Nk} = 0 \quad (\forall k = 1,\dots,N)$.
Insbesondere gilt dann
\begin{align}
\det (T_\text{GS}) &= \det(CC^{-1}) \det(I - CA) = \\
&=  \det(C) \det (C^{-1} - A) = \det(C) \det(M) = 0.
\end{align}
Es muss also mindestens ein Eigenwert von $T_\text{GS}$ gleich Null sein.


\section*{Aufgabe 3}
Gegeben sind die Matrizen
\begin{align}
A_1 &= \begin{pmatrix}
1 & 0 & -2 \\ 0 & 2 & -4 \\ 4 & 2 & 2
\end{pmatrix}&&\text{und}& A_2 &= \begin{pmatrix}
1 & 2 & -2 \\ 1 & 1 & 1 \\ 2 & 2 & 1
\end{pmatrix}
\end{align}
und ein Vektor $b\in\mathbb{R}^3$.
Für das Gauss-Seidel-Verfahren ist $C = (D+L)^{-1}$, wobei $D$ nur die Diagonalelemente und $L$ die Elemente unterhalb der Diagonalen von $A$ enthält.
\subsection*{Für $A_1$:}
Invertieren: 
\begin{align}
D+L = \begin{pmatrix}
1 & 0 & 0 \\ 0 & 2 & 0 \\ 4 & 2 & 2
\end{pmatrix} &\Bigg\vert \begin{pmatrix}
1 & 0 & 0 \\ 0 & 1 & 0 \\ 0 & 0 & 1
\end{pmatrix},\\
\begin{pmatrix}
1 & 0 & 0 \\ 0 & 1 & 0 \\ 0 & 0 & 2
\end{pmatrix} &\Bigg\vert \begin{pmatrix}
1 & 0 & 0 \\ 0 & \frac{1}{2} & 0 \\ -4 & -1 & 1
\end{pmatrix},\\
\begin{pmatrix}
1 & 0 & 0 \\ 0 & 1 & 0 \\ 0 & 0 & 1
\end{pmatrix} &\Bigg\vert \begin{pmatrix}
1 & 0 & 0 \\ 0 & \frac{1}{2} & 0 \\ -2 & -\frac{1}{2} & \frac{1}{2}
\end{pmatrix} = C.
\end{align}
Die Iterationsmatrix $T\equiv T_\text{GS}$ lautet damit
\begin{align}
T &= I - CA = I - \begin{pmatrix}
1 & 0 & 0 \\ 0 & \frac{1}{2} & 0 \\ -2 & -\frac{1}{2} & \frac{1}{2}
\end{pmatrix}\begin{pmatrix}
1 & 0 & -2 \\ 0 & 2 & -4 \\ 4 & 2 & 2
\end{pmatrix} = I - \begin{pmatrix}
1 & 0 & -2 \\ 0 & 1 & -2 \\ 0 & 0 & 7
\end{pmatrix} = \\
&= \begin{pmatrix}
0 & 0 & 2 \\ 0 & 0 & 2 \\ 0 & 0 & -6
\end{pmatrix}.
\end{align}
Wir wollen die Eigenwerte bestimmen, und lösen dafür das Gleichungssystem $Tx = \lambda x$.
Explizit lautet dieses hier
\begin{align}
\begin{pmatrix}
2x_3 \\ 2x_3 \\ -6x_3
\end{pmatrix} &= \lambda\begin{pmatrix}
x_1 \\ x_2 \\ x_3
\end{pmatrix}.
\end{align}
Eine mögliche Lösung ist $\lambda = -6$ mit dem Eigenvektor $x = (1,1,-3)^\top$.
Damit ist der Spektralradius $\rho(T) \le |\lambda| = 6$, insbesondere also größer 1.
Die Iteration konvergiert daher nicht für beliebige Startwerte.

\subsection*{Für $A_2$:}
Invertieren: 
\begin{align}
D+L = \begin{pmatrix}
1 & 0 & 0 \\ 1 & 1 & 0 \\ 2 & 2 & 1
\end{pmatrix} &\Bigg\vert \begin{pmatrix}
1 & 0 & 0 \\ 0 & 1 & 0 \\ 0 & 0 & 1
\end{pmatrix},\\
\begin{pmatrix}
1 & 0 & 0 \\ 0 & 1 & 0 \\ 0 & 2 & 1
\end{pmatrix} &\Bigg\vert \begin{pmatrix}
1 & 0 & 0 \\ -1 & 1 & 0 \\ -2 & 0 & 1
\end{pmatrix},\\
\begin{pmatrix}
1 & 0 & 0 \\ 0 & 1 & 0 \\ 0 & 0 & 1
\end{pmatrix} &\Bigg\vert \begin{pmatrix}
1 & 0 & 0 \\ -1 & 1 & 0 \\ 0 & -2 & 1
\end{pmatrix} = C.
\end{align}
Die Iterationsmatrix $T\equiv T_\text{GS}$ lautet damit
\begin{align}
T &= I - CA = I - \begin{pmatrix}
1 & 0 & 0 \\ -1 & 1 & 0 \\ 0 & -2 & 1
\end{pmatrix}\begin{pmatrix}
1 & 2 & -2 \\ 1 & 1 & 1 \\ 2 & 2 & 1
\end{pmatrix} = I - \begin{pmatrix}
1 & 2 & -2 \\ 0 & -1 & 3 \\ 0 & 0 & -1
\end{pmatrix} = \\
&= \begin{pmatrix}
0 & -2 & 2 \\ 0 & 2 & -3 \\ 0 & 0 & 2
\end{pmatrix}.
\end{align}
Wir wollen die Eigenwerte bestimmen, und lösen dafür das Gleichungssystem $Tx = \lambda x$.
Explizit lautet dieses hier
\begin{align}
\begin{pmatrix}
-2x_2 + 2x_3 \\ 2x_2 - 3x_3 \\ 2x_3
\end{pmatrix} &= \lambda\begin{pmatrix}
x_1 \\ x_2 \\ x_3
\end{pmatrix}.
\end{align}
Eine mögliche Lösung ist $\lambda = 2$ mit dem Eigenvektor $x = (-1,1,0)^\top$.
Damit ist der Spektralradius $\rho(T) \le |\lambda| = 2$, insbesondere also größer 1.
Die Iteration konvergiert daher nicht für beliebige Startwerte.



    



    



\end{document}
